\documentclass[10pt,a4paper]{amsart}
\usepackage[margin=19mm]{geometry}
\usepackage{amsmath,amssymb,mathtools}
\usepackage[scr=rsfs,cal=euler]{mathalfa}
\usepackage{xfrac}
\usepackage[unicode,hidelinks,bookmarks=false]{hyperref}
\newcommand{\ie}{i.e.\ }
\newcommand{\cf}{cf.\ }
\newcommand{\resp}{respectively\ }
\newcommand{\Subsection}{\subsection}
\providecommand{\nobreakdash}{\nobreak\hbox{-}}
\setlength{\emergencystretch}{2em}
\multlinegap0pt
\usepackage[shortlabels]{enumitem}
\usepackage{amssymb}
\usepackage{xcolor}
\usepackage{tikz}
\newtheorem{proposition}{Proposition}
\DeclareMathOperator{\Gen}{Gen}
\newcommand{\codim}{\operatorname{codim}}
\newcommand{\divi}[1]{\langle #1 \rangle}
\DeclareMathOperator{\ev}{ev}
\renewcommand{\geq}{\geqslant}
\renewcommand{\leq}{\leqslant}
\title{High dimensional Riemann--Roch spaces in linear spaces with small squares}
\author{Nihan Tan\i sal\i~, Mieke Wessel}
\date{Source: arXiv:2609.06822v1 (CC BY 4.0)}
\begin{document}
\maketitle

\noindent\emph{Source and licence.} High dimensional Riemann--Roch spaces in linear spaces with small squares. Version arXiv:2609.06822v1 (CC BY 4.0), distributed by arXiv under the Creative Commons Attribution 4.0 International licence, http://creativecommons.org/licenses/by/4.0/, as recorded in the arXiv licence metadata for this exact version and shown on the abstract page https://arxiv.org/abs/2609.06822v1. Full licence notice: \texttt{licenses/arxiv-2609.06822-CC-BY-4.0.txt}.\par\smallskip

\noindent\textit{Editorial scope.} Counted unit: the article's proposition prop:m\_dimension (source0.tex lines 615--623). The article's complete original proof of it is delivered with it (source0.tex lines 624--702, one \texttt{proof} environment of 79 lines). No mathematics is written, repaired or re-derived: the statement and the proof are the article's own lines, in the article's own notation, and no retained line is substituted or re-flowed.  No further source-local context is needed: every cross-reference of every retained line resolves inside the two delivered spans.  Nothing was omitted from any retained span, so the package carries no omission record.  No retained line contains a citation, so the wrapper carries no bibliography and no bibliography entry.  No retained line contains a figure environment, an image-inclusion command or drawing code, so no image is delivered and the package declares no asset dependency.  Editorial formatting only, in addition: an AMS \texttt{amsart} preamble that restates the article's own declarations and macros used by the retained lines, namely the environments \texttt{proposition} on separate counters, together with the standard packages those lines need.  The retained environments print in this document as Proposition 1 in order of appearance, whereas the article prints the same items with the numbering of its own full document.  The complete native source of the article as distributed in the arXiv e-print for this exact version is bundled unchanged as \texttt{sources/209575/source0.tex}.
	\begin{proposition}\label{prop:m_dimension}
		Suppose that $U \subseteq L(A)$ is a $K$-subspace, so that
		$U^2 \subseteq L(2A)$.
		Let $\varphi \in \Phi_{2A }(U^2)$ be a form and $X$ an expression of weight $w$ such that $\varphi= X_{2A}$.
		Define
		$$T:=U\cap\bigcap_{\chi\in \divi{D_X} _A}\ker(\chi).$$
		Then we have
		$$\operatorname{codim}_U(T) \le \frac{w}{2}.$$
	\end{proposition}

	\begin{proof}
		When $\dim (U) \leq \frac w 2$ the statement is trivial, so we assume that $\dim U > \frac w2$. The rest of the lemma we prove by contradiction, so suppose that $\codim_U(T) = x > w/2$. 
		
		Write $\operatorname{Gen}_A(X)$ for the set of generating evaluation forms $\ev_A(P, j)$ of $\divi {D_X}_A$, that is
		$$ \Gen_{A}(X)= \{\ev_{A}(P_i,j):1\le i\le \sigma,\ 0\le j\le r_i\} $$
		where $r_i= r_{P_i} (X)$. We can then equivalently define $T = U \cap \bigcap_{\chi \in \Gen_{A}(X)} \ker(\chi).$
		
		We order $\Gen_A(X)$ lexicographically: first by the index $i$ of the places and then by the index $j$, both in increasing order. This order is denoted by $\leq$. For example, we have $\ev_A(P_1, 1) \leq \ev(P_1, 3) \leq \ev(P_2, 0).$ 
		
		Choose elements 
		$$\psi_1, \ldots , \psi_x \in \Gen_A(X)  $$
		inductively as follows
		\begin{enumerate}[1.]
			\item Set $U_0:=U$.
			\item For each $1\le \ell \le x$, let $\psi_\ell$ be the smallest element of $\Gen_A(X)$ with respect to $\leq$ such that
			$$
			U_\ell:=U \cap\bigcap_{\chi\in\Gen_A(X), \chi\le \psi_\ell }\ker(\chi)
			\subsetneq U_{\ell-1}
			$$
			and $\operatorname{codim}_U(U_\ell )=\ell$.
		\end{enumerate}
		For $\ell = 1, \ldots, x$, define $i_\ell$ and $j_\ell$ such that
		$$\psi_{\ell}=\ev_{A}(P_{i_\ell},j_\ell).$$
		Then $T = U \cap \bigcap_{\ell = 1}^x \ker(\psi_\ell)$.
		We choose $u_\ell \in U_{\ell-1} \setminus U_\ell $ so that 
		\begin{equation} \label{eq:ul}
			\ev_A(P_i,j)(u_\ell) =\begin{cases}
				1&\text{if } i=i_\ell \qquad \text{and} \qquad j=j_\ell;\\
				0&\text{if } i= i_\ell \qquad\text{and}\qquad  j< j_\ell;\\
				0&\text{if } i< i_\ell.
			\end{cases}
		\end{equation} 
		Then in particular we have that
		\begin{align}
			v_{P_{i_\ell} } (u_\ell) =-v_{P_{i_\ell}} (A) +j_\ell .
			\label{eq:s_u}
		\end{align}
		Now, we define 
		\begin{alignat*}{2}
			& \Gamma_1 := \{\psi_1, \ldots, \psi_x\}, &&\\
			&\Gamma_2:=\{\ev_A( P, j) \in \Gen_A(X) \mid \ev_A (P,r_P(X) - j)\in \Gamma_1\}, &&\\
			& B:= \Gen_A(X) \setminus \Gamma_2 \quad \text{ and }&&\\
			& V := U \cap \bigcap_{\chi\in B}\ker(\chi). &&
		\end{alignat*}
		Then $|B|=w-x<x $ and thus, $\dim V >\dim T.$ Let $u_0 \in V\setminus T$ and note that $T = V \cap \bigcap_{\chi \in \Gamma_2} \ker(\chi)$. Because $u_0\notin T$, there must be some element $\psi_0 \in \Gamma_2$ such that $\psi_0(u_0) \neq 0$. Hence, we can take $i_0$ maximal such that for some $j$ we have that
		$$\ev_A(P_{i_0},j )\in \Gamma_2 \qquad\text{and}\qquad \ev_A(P_{i_0},j)(u_0)\neq 0.$$
		We define $j_0$ to be the minimal $j$ for which this holds. Then, in particular, we have 
		\begin{align}
			v_{P_{ i_0} } (u_0)  = -v_{P_{i_0}}(A) +j_0    \label{eq:d}.
		\end{align}
		Note that by the construction of $\Gamma_2$, there exists $\ell$ such that 
		\begin{equation}\label{eq:evpsi} \psi_\ell=\ev_{A}(P_{i_0},r_{i_0}-j_0) \in \Gamma_1.
		\end{equation}
		From now on we fix $\ell$ such that \eqref{eq:evpsi} holds, that is $i_\ell = i_0$ and $j_\ell = r_{i_0} - j_0$. 
		
		After possibly normalizing $u_0$, we claim that for $i \in [1, \sigma]$ and $j \in [0, r_i]$ we have
		$$ \ev_{ 2A}(P_i,j) (u_0 u_\ell) = \begin{cases}
			1 &\quad \text{if}\quad i=i_0 \quad\text{and}\quad j=r_{i_0};\\0& \quad\text{otherwise}.
		\end{cases}$$
		
		When $i \neq i_0$ we can use the maximality of $i_0$ and Equation \eqref{eq:ul} to find
		\begin{align*}\ev_A(P_i,j) (u_0) =0 
			&\quad\text{if}\qquad i> i_0 \qquad \text{ and } \qquad  j\le r_i,\\
			\ev_A(P_i,j) (u_\ell) =0 &\quad\text{if}\qquad i< i_0 \qquad \text{ and } \qquad j\le r_i.
		\end{align*}
		This implies that $v_{P_i}(u_0 u_\ell) \geq -v_{P_i}(2A) + r_i + 1$ for these $i$ and thus, the claim for $i \neq i_0$ must hold.
		
		When $i = i_0$, by Equations \eqref{eq:s_u} and \eqref{eq:d}, 
		we have 
		$$ v_{P_{i_0}}(u_0 u_\ell) = -v_{P_{i_0}}(2A)+ r_{i_0}.  $$ 
		This directly implies the claim for $j < r_{i_0}$ and because we are allowed to normalize $u_0$ we also get the result for $j = r_{i_0}$. 
		
		Finally, recall that $\varphi$ has the following decomposition
		$$\varphi = \lambda \ev_{2A} (P_{i_0} ,  r_{i_0}) + \sum_{ (i,j)\neq (i_0,r_{i_0})}\lambda_{i j} \ev_{2A}(P_i,j ), $$
		for some non-zero $\lambda$. Hence, evaluating $\varphi$ at
		$u_0 u_\ell \in U^2$ we get
		$$\varphi (u_0 u_\ell) = \lambda + 0\neq 0. $$
		This contradicts the fact that $\varphi\in \Phi_{2A}(U^2).$ We conclude that $x \leq \frac w2$, as wished.
	\end{proof}

\end{document}
